\chapter*{使用说明：本册是怎么筛题的}
\addcontentsline{toc}{chapter}{使用说明：本册是怎么筛题的}
\markboth{使用说明：本册是怎么筛题的}{}

\section*{一、本册来源与配套关系}

本册题目全部选自\textbf{《工程流体力学 学习指导及习题解答》}
（\textbf{陈洁、袁铁江 编著}，清华大学出版社，2015 年 2 月第 1 版，ISBN 978-7-302-37859-4）。
该书是\textbf{孔珑 主编《工程流体力学（第四版）》}的配套辅导书，体例与教材相同（共十章，
每章由"主要内容 / 本章难点 / 课后习题解答"三部分组成），书末\textbf{只给出课后习题的详细解答}。
孔珑教材是全国"十一五""十二五"国家级规划教材，题目质量高、覆盖全，
是赵琴教材课后习题之外的\textbf{最优质补充题源}。

\noindent\textbf{为什么可以拿孔珑的题练 818？}因为 818 考纲（赵琴《流体力学与流体机械》第 1--8 章）
所覆盖的\textbf{流体静力学、流体运动学、伯努利方程、动量方程、管路水力计算、相似与量纲、
势流与旋涡}等核心内容，两本书的物理体系完全一致，题型与解法直接通用。

\section*{二、筛选原则：只留符合 818 考纲的题}

\begin{center}
\begin{tabularx}{\textwidth}{@{}lX@{}}
\toprule
孔珑书中内容 & 处理方式 \\
\midrule
流体静力学、流体运动学、流体动力学基础、管路水力计算、相似理论与量纲分析、势流与旋涡 & \textbf{全部收录} \\
明渠流、堰流、渗流、边界层与绕流（超出部分） & 只收与赵琴第 5、7、8 章对应的基础题 \\
气体动力学（一元可压缩流动） & \textbf{不收}（赵琴第 9 章，初试不考） \\
流体机械（泵、风机、水轮机） & \textbf{不收}（赵琴第 11--14 章，复试范围） \\
\bottomrule
\end{tabularx}
\end{center}

\section*{三、章号对照（孔珑书 $\leftrightarrow$ 赵琴书）}

\begin{center}
\begin{tabularx}{\textwidth}{@{}clX@{}}
\toprule
本册章 & 对应赵琴章 & 内容 \\
\midrule
第 1 章 & 赵琴 1.2--1.4 & 流体的物理性质 \\
第 2 章 & 赵琴 2 & 流体静力学 \\
第 3 章 & 赵琴 3 & 流体运动学 \\
第 4 章 & 赵琴 4 & 流体动力学基础（伯努利、动量方程） \\
第 5 章 & 赵琴 5 & 管路、孔口、管嘴水力计算 \\
第 6 章 & 赵琴 6 & 相似理论与量纲分析 \\
第 7 章 & 赵琴 7 & 势流与旋涡 \\
第 8 章 & 赵琴 8 & 黏性流体动力学基础（N-S、边界层） \\
\bottomrule
\end{tabularx}
\end{center}

\section*{四、星级怎么定的}

与《题库汇编》同一套标准，\textbf{只由文件证据决定，不凭印象}：

\begin{center}
\begin{tabularx}{\textwidth}{@{}clX@{}}
\toprule
星级 & 含义 & 判定条件（满足其一） \\
\midrule
\xstarfull{5} & 必做 & 对应考点在历年真题中作为计算题出现 $\ge 3$ 次，或考情分析列为"考试计算题重点" \\
\xstarfull{4} & 高频 & 真题中出现 2--3 次，或考情分析写"重点掌握" \\
\xstarfull{3} & 常考 & 真题中出现 1 次，或考点分析写"掌握基本概念与公式套用" \\
\xstarfull{2} & 了解 & 真题未出现，但属考纲核心内容 \\
\xstarfull{1} & 边缘 & 仅作扩展 \\
\bottomrule
\end{tabularx}
\end{center}

\noindent 题目在前、解答统一置于书末，方便先自测再对答案。
